% Préambule commun des schémas TikZ — charte Liouaï
% Chaque schéma commence par % Préambule commun des schémas TikZ — charte Liouaï
% Chaque schéma commence par % Préambule commun des schémas TikZ — charte Liouaï
% Chaque schéma commence par \input{_preambule} puis \begin{document}.
\documentclass[tikz,border=6pt]{standalone}
\usepackage[T1]{fontenc}
\usepackage[utf8]{inputenc}
\usepackage[scaled=0.95]{helvet}
\renewcommand{\familydefault}{\sfdefault}
\usepackage{amsmath,amssymb}
\usepackage[european]{circuitikz}
\usetikzlibrary{arrows.meta,positioning,calc,angles,quotes,decorations.pathreplacing,backgrounds}

% Couleurs du logo Liouaï Digital Education
\definecolor{navy}{HTML}{243743}
\definecolor{vert}{HTML}{28B78D}
\definecolor{vertfonce}{HTML}{167A5C}
\definecolor{vertclair}{HTML}{EAF7F2}
\definecolor{gris}{HTML}{A1A9AE}
\definecolor{gristexte}{HTML}{5C6870}
\definecolor{fond}{HTML}{F3F5F6}
\definecolor{jaune}{HTML}{E5C800}
\definecolor{jauneclair}{HTML}{FFFBD6}

\tikzset{
  every picture/.style={line cap=round, line join=round, color=navy, text=navy},
  fleche/.style={-{Stealth[length=2.6mm,width=2mm]}, line width=1.1pt},
  vecteur/.style={fleche, line width=1.4pt},
  titre/.style={font=\bfseries\small, text=navy},
  note/.style={font=\footnotesize, text=gristexte},
  boite/.style={draw=navy, line width=0.9pt, rounded corners=5pt, fill=fond, align=center, inner sep=6pt},
}
 puis \begin{document}.
\documentclass[tikz,border=6pt]{standalone}
\usepackage[T1]{fontenc}
\usepackage[utf8]{inputenc}
\usepackage[scaled=0.95]{helvet}
\renewcommand{\familydefault}{\sfdefault}
\usepackage{amsmath,amssymb}
\usepackage[european]{circuitikz}
\usetikzlibrary{arrows.meta,positioning,calc,angles,quotes,decorations.pathreplacing,backgrounds}

% Couleurs du logo Liouaï Digital Education
\definecolor{navy}{HTML}{243743}
\definecolor{vert}{HTML}{28B78D}
\definecolor{vertfonce}{HTML}{167A5C}
\definecolor{vertclair}{HTML}{EAF7F2}
\definecolor{gris}{HTML}{A1A9AE}
\definecolor{gristexte}{HTML}{5C6870}
\definecolor{fond}{HTML}{F3F5F6}
\definecolor{jaune}{HTML}{E5C800}
\definecolor{jauneclair}{HTML}{FFFBD6}

\tikzset{
  every picture/.style={line cap=round, line join=round, color=navy, text=navy},
  fleche/.style={-{Stealth[length=2.6mm,width=2mm]}, line width=1.1pt},
  vecteur/.style={fleche, line width=1.4pt},
  titre/.style={font=\bfseries\small, text=navy},
  note/.style={font=\footnotesize, text=gristexte},
  boite/.style={draw=navy, line width=0.9pt, rounded corners=5pt, fill=fond, align=center, inner sep=6pt},
}
 puis \begin{document}.
\documentclass[tikz,border=6pt]{standalone}
\usepackage[T1]{fontenc}
\usepackage[utf8]{inputenc}
\usepackage[scaled=0.95]{helvet}
\renewcommand{\familydefault}{\sfdefault}
\usepackage{amsmath,amssymb}
\usepackage[european]{circuitikz}
\usetikzlibrary{arrows.meta,positioning,calc,angles,quotes,decorations.pathreplacing,backgrounds}

% Couleurs du logo Liouaï Digital Education
\definecolor{navy}{HTML}{243743}
\definecolor{vert}{HTML}{28B78D}
\definecolor{vertfonce}{HTML}{167A5C}
\definecolor{vertclair}{HTML}{EAF7F2}
\definecolor{gris}{HTML}{A1A9AE}
\definecolor{gristexte}{HTML}{5C6870}
\definecolor{fond}{HTML}{F3F5F6}
\definecolor{jaune}{HTML}{E5C800}
\definecolor{jauneclair}{HTML}{FFFBD6}

\tikzset{
  every picture/.style={line cap=round, line join=round, color=navy, text=navy},
  fleche/.style={-{Stealth[length=2.6mm,width=2mm]}, line width=1.1pt},
  vecteur/.style={fleche, line width=1.4pt},
  titre/.style={font=\bfseries\small, text=navy},
  note/.style={font=\footnotesize, text=gristexte},
  boite/.style={draw=navy, line width=0.9pt, rounded corners=5pt, fill=fond, align=center, inner sep=6pt},
}

\begin{document}
\begin{tikzpicture}[scale=0.85]
\newcommand{\axes}{%
  \draw[fleche, gris, line width=0.7pt] (0,0) -- (3.6,0);
  \draw[fleche, gris, line width=0.7pt] (0,0) -- (0,3);}
% --- Une solution
\begin{scope}
\axes
\draw[vertfonce, line width=1.6pt] (0.2,0.4) -- (3.3,2.6);
\draw[navy, line width=1.6pt] (0.2,2.5) -- (3.3,0.5);
\fill[jaune, draw=navy] (1.67,1.44) circle (3pt);
\node[titre] at (1.8,3.6) {une solution};
\node[note, align=center] at (1.8,-0.75) {droites sécantes};
\node[font=\footnotesize, align=center] at (1.8,-1.55) {$\operatorname{rg}A=\operatorname{rg}[A\,|\,\mathbf{b}]=n$};
\end{scope}
% --- Aucune solution
\begin{scope}[xshift=5.2cm]
\axes
\draw[vertfonce, line width=1.6pt] (0.2,0.3) -- (3.3,1.9);
\draw[navy, line width=1.6pt] (0.2,1.1) -- (3.3,2.7);
\node[titre] at (1.8,3.6) {aucune solution};
\node[note, align=center] at (1.8,-0.75) {droites parallèles distinctes};
\node[font=\footnotesize, align=center] at (1.8,-1.55) {$\operatorname{rg}A<\operatorname{rg}[A\,|\,\mathbf{b}]$};
\end{scope}
% --- Une infinité
\begin{scope}[xshift=10.4cm]
\axes
\draw[vertfonce, line width=4pt, opacity=0.45] (0.2,0.5) -- (3.3,2.4);
\draw[navy, line width=1.4pt, densely dashed] (0.2,0.5) -- (3.3,2.4);
\node[titre] at (1.8,3.6) {une infinité};
\node[note, align=center] at (1.8,-0.75) {droites confondues};
\node[font=\footnotesize, align=center] at (1.8,-1.55) {$\operatorname{rg}A=\operatorname{rg}[A\,|\,\mathbf{b}]<n$};
\end{scope}
\end{tikzpicture}
\end{document}
